\section{Introduction}

AI systems increasingly generate both solutions and the artifacts used to evaluate those solutions.
In formal verification workflows, these artifacts include tests, theorem statements, and proof-oriented specifications in Lean~4.
This creates a measurement bottleneck: humans can audit only a small fraction of generated artifacts, yet deployment decisions depend on reliable evaluation.

Current LLM-as-a-judge practice often treats human agreement as the main validity target.
That target is necessary but insufficient.
A judge may correlate with humans on a finite panel while violating basic behavioral constraints that any trustworthy evaluator should satisfy.
For example, it may fail to assign maximal scores to identical inputs, fail to penalize progressively worse artifacts, or produce unstable scores under small prompt changes.
If these failures are not detected, benchmark rankings can be brittle or strategically gamed.

We frame this as an AI measurement science problem.
A valid judge should satisfy explicit, falsifiable properties under controlled perturbations.
In this paper, we operationalize that framing with \textsc{CertJudge}, a certification protocol for Lean~4 judge reliability.
The protocol tests identity (P1), monotonicity to injected bugs (P2), monotonicity to specification ablations (P3), and run-level stability (P4), and then aggregates them into a trust index for judge selection.

This framing aligns with three core problems in AI Measurement Science.
First, it supports \emph{interactive measurement} by evaluating judges in agentic settings where artifacts are generated end-to-end rather than in isolated subtasks.
Second, it addresses \emph{measurement under strategic optimization} by making important failure modes auditable before judge deployment.
Third, it supports \emph{measurement under non-stationarity} by enabling repeated recalibration as prompts, models, and task mixtures evolve.

Beyond this specific Lean~4 setting, the key claim is measurement-theoretic:
evaluation quality must be measured and governed as a first-class object.
Property certification supplies a concrete validity layer that can be reused across domains where evaluators are expensive, gameable, and subject to drift.

Our contributions are:
\begin{enumerate}
\item \textbf{A measurement-science framing for LLM judges in formal verification.}
We recast judge validation from one-shot correlation to property-based certification with explicit validity constraints.
\item \textbf{A concrete certification protocol for Lean~4 judge reliability with general measurement value.}
We provide a reproducible four-property suite and trust-index aggregation for amortized judge screening.
\item \textbf{Empirical evidence on calibration and downstream ranking.}
On a human-labeled validation panel, the trust index tracks judge-level human alignment, and on \textsc{VeriBench} the certified judge recovers meaningful method ordering between baseline and feedback-driven agents.
\end{enumerate}
